% p.70 — العوامل ومتعلقاتها (تتمة)
% Layout: twocol, as the lithograph divides the page into two columns of numbered entries
% (number + governing word | example phrases) with column-width rule paragraphs (منها …)
% between the groups. rightcol: entries 50–53, rule, 54–57, rule, 58–64; leftcol: 65–66, rule,
% 67–79, rule (افعال المقاربة). Rule paragraphs are \par-separated text between the sarf tables;
% entries with several example lines are extra rows with an empty first cell.
% Note: reledpar sets the two columns line by line (line i of the right column beside line i of
% the left one, each pair as tall as its taller member), so a column made of tall sarf "lines"
% padded the other column and the page overflowed. Each column's whole content is therefore
% wrapped in one rightbox/leftbox (minipages inside the reledpar columns): the columns then
% flow independently and the page fits; the apparatus numbers each column as a single line.
\mspage{70}
\begin{twocol}
\begin{rightbox}
\begin{sarf}{|c|c|}
(50) اَيْنَ & اَيْنَ تَجْلِسْ اَجْلِسْ\\
(51) اَنّٰى & اَنّٰى تُعَامِلْ اُعَامِلْ\\
(52) حَيْثُمَا & حَيْثُمَا تَأْكُلْ آكُلْ\\
(53) اِذْمَا & اِذْمَا تَنْزِلْ اَنْزِلْ\\
\end{sarf}
مِنْهَا اَسْمَاءٌ تَنْصِبُ الْاِسْمَ عَلَى التَّمْيِيْزِ وَهِيَ اَرْبَعَةٌ
\par
\begin{sarf}{|c|c|}
(54) عَشَرَ & رَاَيْتُ اَحَدَ عَشَرَ كَوْكَبًا\\
(55) كَمْ & كَمْ دِيْنَارًا اَعْطَيْتَ زَيْدًا\\
 & كَمْ عِنْدِيْ دِيْنَارًا\\
(56) كَاَيِّنْ & كَاَيِّنْ رَجُلًا لَقِيْتُ\\
 & كَاَيِّنْ رَجُلًا عِنْدَكَ\\
(57) كَذَا & عِنْدِيْ كَذَا كَذَا مَالًا\\
\end{sarf}
مِنْهَا كَلِمَاتٌ تُسَمّٰى اَسْمَاءَ الْاَفْعَالِ بَعْضُهَا تَرْفَعُ الْفَاعِلَ وَهِيَ ثَلَاثَةٌ وَبَعْضُهَا تَنْصِبُ الْمَفْعُوْلَ وَهِيَ سِتَّةٌ
\par
\begin{sarf}{|c|c|}
(58) رُوَيْدَ & رُوَيْدَ زَيْدًا اَيْ اَمْهِلْ\\
(59) بَلْهَ & بَلْهَ بَكْرًا اَيْ اُتْرُكْ\\
(60) دُوْنَكَ & دُوْنَكَ خَالِدًا اَيْ \com{اُتْرُكْ}{كذا في الأصل، ولعل الصواب: خُذْ}\\
(61) عَلَيْكَ & عَلَيْكَ رَحِيْمًا اَيْ \unclear{اَلْزِمْ}\\ % TODO: vowels of اَلْزِمْ (expected اِلْزَمْ)
(62) \com{هَيَّهَلْ}{كذا في الأصل، والحرف الأول غير واضح (ه/ح)، ولعل الصواب: حَيَّهَلْ} & \unclear{هَيَّهَلْ} هَادِيًا اَيْ اِئْتِ بِهِ\\
(63) هَا & هَا \unclear{دَلُوًا} اَيْ خُذْ\\ % TODO: second word unclear; perhaps دَلْوًا
(64) هَيْهَاتَ & هَيْهَاتَ الْاَمْرُ اَيْ بَعُدَ\\
\end{sarf}
\end{rightbox}
\begin{leftbox}
\begin{sarf}{|c|c|}
(65) شَتَّانَ & شَتَّانَ زَيْدٌ وَعَمْرٌو اَيْ اِفْتَرَقَ\\
(66) سَرْعَانَ & سَرْعَانَ خَالِدٌ اَيْ اَسْرَعَ\\
\end{sarf}
مِنْهَا اَفْعَالُ النَّاقِصَةِ تَرْفَعُ الْاِسْمَ وَتَنْصِبُ الْخَبَرَ وَهِيَ ثَلَاثَةَ عَشَرَ فِعْلًا
\par
\begin{sarf}{|c|c|}
(67) كَانَ & كَانَ زَيْدٌ قَائِمًا كَانَ الْفَقِيْرُ غَنِيًّا\\
 & كَانَ اللهُ عَلِيْمًا حَكِيْمًا كَانَ الْمَطَرُ\\
(68) صَارَ & صَارَ زَيْدٌ قَائِمًا صَارَ الْجَاهِلُ عَالِمًا\\
(69) اَصْبَحَ & اَصْبَحَ زَيْدٌ فَقِيْهًا اَصْبَحَ زَيْدٌ مَرِيْضًا\\
(70) اَمْسٰى & اَمْسٰى زَيْدٌ عَالِمًا اَمْسٰى زَيْدٌ صَحِيْحًا\\
(71) اَضْحٰى & اَضْحٰى زَيْدٌ حَاكِمًا اَضْحٰى زَيْدٌ نَاصِرًا\\
(72) ظَلَّ & ظَلَّ زَيْدٌ كَاتِبًا ظَلَّ الصَّبِيُّ بَالِغًا\\
(73) بَاتَ & بَاتَ الشَّبَابُ شَيْخًا بَاتَ زَيْدٌ نَائِمًا\\
(74) مَازَالَ & مَازَالَ اللهُ حَيًّا قَيُّوْمًا\\
(75) مَابَرِحَ & مَابَرِحَ زَيْدٌ صَائِمًا\\
(76) مَافَتِئَ & مَافَتِئَ عَبْدُ اللهِ مُصَلِّيًا\\
(77) مَاانْفَكَّ & مَاانْفَكَّ عَبْدُ اللهِ حَلِيْمًا\\
(78) مَادَامَ & اِجْلِسْ مَادَامَ زَيْدٌ جَالِسًا\\
(79) لَيْسَ & لَيْسَ زَيْدٌ قَائِمًا\\
 & لَيْسَ زَيْدٌ بِقَائِمٍ\\
\end{sarf}
مِنْهَا اَفْعَالٌ تُسَمّٰى اَفْعَالَ الْمُقَارَبَةِ تَرْفَعُ الْفَاعِلَ وَخَبَرُهُ يَكُوْنُ مُضَارِعًا مَعَ اَنْ اَوْ بِلَا اَنْ
\par
\end{leftbox}
\end{twocol}
